We see that in the Thomas--Fermi model the charge-density distributions in different atoms are similar, with the charac%
\SourcePageJoin{322}{325}
teristic length scale set by $Z^{-1/3}$ (in ordinary units: $\hbar^2/(me^2Z^{1/3})$, i.e. the Bohr radius divided by $Z^{1/3}$). If distances are measured in atomic units, then, in particular, the distances at which the electron density is greatest will be the same for all $Z$. It may therefore be stated that most of the electrons in an atom of atomic number $Z$ are located at distances from the nucleus of order $Z^{-1/3}$. Calculation shows that half the atom's total electronic charge lies inside a sphere of radius $1{,}33Z^{-1/3}$.
